%! Characteristics of Exponential Functions — Unit 6, Lesson 6.0 — Student Packet Cover Sheet
\documentclass[10pt]{article}
\usepackage{algebra2-article}
\usepackage{algebra2-boxes}
\usepackage{ltablex}
\keepXColumns

\begin{document}

% Full-bleed forest banner
\begin{tikzpicture}[remember picture, overlay]
  \fill[forest] (current page.north west)
    rectangle ([yshift=-0.9in]current page.north east);
\end{tikzpicture}
\vspace{-0.8in}

\begin{center}
  {\color{white}\textbf{\LARGE Algebra 2: Shepherd}}\\[4pt]
  {\color{forestmid}\large\textbf{Unit 6 \quad Exponential Functions}}\\[6pt]
  {\color{white}\textbf{\large Lesson 6.0 \quad Characteristics of Exponential Functions}}
\end{center}

\vspace{0.22in}
\namedateperiod
\vspace{0.18in}

\begin{learningtargetbox}
\small
\begin{enumerate}[leftmargin=1.5em, itemsep=5pt]
  \item \ldots{} identify an \textbf{exponential} function from its equation, its graph, or a \textbf{table}, using the \textbf{constant ratio} fingerprint --- and tell \textbf{growth} from \textbf{decay} by comparing the base $b$ to $1$.
  \item \ldots{} read the \textbf{domain}, \textbf{range}, \textbf{intercepts}, \textbf{increasing/decreasing} and \textbf{positive/negative} intervals, \textbf{extrema}, \textbf{horizontal asymptote}, and \textbf{end behavior} off an exponential graph.
  \item \ldots{} explain why an exponential graph has \textbf{no} $x$-intercept and \textbf{no} absolute minimum, even though it clearly flattens out along a floor --- and \textbf{compare and contrast} exponential functions with the families from Units 1--5.
\end{enumerate}
\end{learningtargetbox}

\vspace{0.14in}

\begin{tocbox}
\small
\renewcommand{\arraystretch}{1.5}
\begin{tabularx}{\linewidth}{@{} c l X r @{}}
\rowcolor{forestbg}
\textbf{\#} & \textbf{Component} & \textbf{Description} & \textbf{Score} \\
\midrule
1 & Warm-Up        & Spiral review: zero and negative exponents, add-versus-multiply tables, and a doubling rumor & \blank{1.2cm} \\
2 & Guided Notes   & Reading every characteristic off the growth and decay parents & \blank{1.2cm} \\
3 & Group Activity & \emph{Reading Exponential Graphs} --- feature tables, features from equations, and a medication model, by tier & \blank{1.2cm} \\
4 & Exit Ticket    & Quick check: features from a graph, plus an SOL-style range item & \blank{1.2cm} \\
5 & Homework       & Features from equations, two contrasting graphs, the table fingerprint, and a depreciation model & \blank{1.2cm} \\
\midrule
  & \multicolumn{2}{r}{\textbf{Total}} & \blank{1.2cm} \\
\end{tabularx}
\end{tocbox}

\vspace{0.10in}

\begin{remindbox}
\small An \textbf{exponential function} puts the variable in the \textbf{exponent}:
$f(x)=ab^{x}$, with $a\ne 0$, $b>0$, and $b\ne 1$. The \textbf{initial value} $a$ is the
$y$-intercept, because $b^{0}=1$ for every legal base. The \textbf{base} $b$ is what you
\textbf{multiply} by each time $x$ goes up by $1$ --- so an exponential table has a \textbf{constant
ratio}, just as a linear table has a constant difference. Compare $b$ to $1$, not to $0$: $b>1$ is
\textbf{growth} and $0<b<1$ is \textbf{decay}. Every graph in this family is one smooth, unbroken,
never-turning curve, so there are \emph{no} turning points and \emph{no} extrema. And because a
positive base raised to any real power is always positive, $b^{x}$ is never $0$ --- so the graph
never crosses the $x$-axis. It only \emph{approaches} the line $y=0$, its \textbf{horizontal
asymptote}. That line is a boundary the range excludes, not a value the function reaches, which is
why the range is $(0,\infty)$ and \emph{not} $[0,\infty)$, and why there is no absolute minimum ---
the opposite of Unit 5, where the square root's \textbf{endpoint} was a point the graph actually
landed on. Unit 5 took away half the \textbf{domain}; Unit 6 hands it back and takes half the
\textbf{range} instead.
\end{remindbox}

\end{document}
